\documentclass[11pt,a4paper]{article}

\usepackage{kotex}
\usepackage{fontspec}
\setmainfont{Times New Roman}
\setmainhangulfont{AppleMyungjo}
\setsanshangulfont{Apple SD Gothic Neo}

\usepackage[a4paper,top=18mm,bottom=18mm,left=20mm,right=20mm]{geometry}
\usepackage{fancyhdr}
\setlength{\headheight}{24pt}
\usepackage{enumitem}
\usepackage{mdframed}
\usepackage{array}
\usepackage{booktabs}

\setlength{\parindent}{1em}
\linespread{1.18}

\pagestyle{fancy}
\fancyhf{}
\renewcommand{\headrulewidth}{0.6pt}
\fancyhead[L]{\sffamily\bfseries 2026 고1 영어 변형 — 지문 40}
\fancyhead[R]{\sffamily 무의식적 위험 감지}
\renewcommand{\footrulewidth}{0pt}
\fancyfoot[C]{\framebox[16mm][c]{\thepage}}

\newcommand{\qno}[1]{\par\vspace{8pt}\noindent{\sffamily\bfseries #1.}\ }
\newcommand{\instr}[1]{{\sffamily #1}\par\vspace{3pt}}
\newlist{choices}{enumerate}{1}
\setlist[choices]{label=,leftmargin=1.5em,itemsep=2pt,topsep=3pt,parsep=0pt}
\newmdenv[linewidth=0.6pt,roundcorner=3pt,innertopmargin=7pt,innerbottommargin=7pt,
          innerleftmargin=9pt,innerrightmargin=9pt,skipabove=5pt,skipbelow=5pt]{pbox}
\newcommand{\nemo}[1]{\fbox{\,#1\,}}

\begin{document}

\begin{center}
{\fontsize{15}{17}\selectfont\sffamily\bfseries [지문 40] 다음 글을 읽고 물음에 답하시오.}
\end{center}
\vspace{4pt}

% ===== 공통 지문 박스 (빈칸: "generate reasons for why it was accurate" → blank) =====
\begin{pbox}
We go out in the world every day and make decisions about what is safe or not and what is appropriate or not. It is what psychologist Joseph LeDoux has suggested is an unconscious ``danger detector'' that determines whether or not something or someone is safe before we can even begin to consciously make a determination. When the object, animal, or person is assessed to be dangerous, a fear response, which has been called ``fight or flight,'' occurs. On a conscious level we may correct a mistake in this ``danger detector'' when we notice it, but often, of course, we simply begin to \rule{4.5cm}{0.4pt} the danger detector's initial response. We are generally convinced that our decisions are ``rational,'' but in reality, most human decisions are made emotionally, and then we collect or generate the facts to justify these decisions. When we see something or someone who ``feels'' dangerous, we have already launched into action internally before we have even started ``thinking.''
\end{pbox}

\vspace{6pt}

% ===== 요약문 (이탤릭, (A)(B) 빈칸) =====
\noindent\hspace{1em}\textit{Our brains make (A)\,\rule{2.8cm}{0.4pt}\, judgments about danger before rational thought kicks in, and we then tend to (B)\,\rule{2.8cm}{0.4pt}\, those judgments with consciously generated reasoning.}

\vspace{10pt}

% ===================== Q1: 요약문 (A)(B) =====================
\qno{1}\instr{다음은 윗글을 요약한 것이다. 빈칸 (A), (B)에 들어갈 말로 가장 적절한 것은?}

\begin{center}
\renewcommand{\arraystretch}{1.25}
\begin{tabular}{c >{\centering\arraybackslash}p{3.8cm} >{\centering\arraybackslash}p{3.8cm}}
 & (A) & (B) \\
① & deliberate & challenge \\
② & preconscious & suppress \\
③ & preconscious & rationalize \\
④ & systematic & validate \\
⑤ & instinctive & ignore \\
\end{tabular}
\end{center}

\vspace{10pt}

% ===================== Q2: 빈칸추론 =====================
\qno{2}\instr{다음 빈칸에 들어갈 말로 가장 적절한 것은?}
\begin{choices}
\item ① challenge
\item ② ignore
\item ③ override
\item ④ justify
\item ⑤ abandon
\end{choices}

\vspace{10pt}

% ===================== Q3: 서술형 해석 =====================
\qno{3}\instr{[서술형] 다음 문장을 우리말로 해석하시오.}

\vspace{4pt}
\noindent When we see something or someone who ``feels'' dangerous, we have already launched into action internally before we have even started ``thinking.''

\vspace{6pt}
\noindent $\Rightarrow$ \rule{\linewidth}{0.4pt}

\vspace{4pt}
\noindent \rule{\linewidth}{0.4pt}

% ===================== Q4 =====================
\qno{4}\instr{윗글의 제목으로 가장 적절한 것은?}
\begin{choices}
\item ① The Rational Mind That Always Corrects Emotional Errors
\item ② How Conscious Thinking Prevents Every Fear Response
\item ③ The Unconscious Detector Behind Emotion-First Decisions
\item ④ Why People Should Ignore Feelings of Danger
\item ⑤ The Complete Separation of Emotion and Decision Making
\end{choices}

\end{document}
